\section{Rousseau’s Third Solitary Walk}

\begin{quotex}
``Growing older, I learn all the time.'' Solon often repeated this line in his old age. In a sense I could say the same, but the knowledge that the experience of twenty years has brought me is a poor thing, and even ignorance would be preferable. No doubt adversity is a great teacher, but its lessons are dearly bought, and often the profit we gain from them is not worth the price they cost us. What is more, these lessons come so late in the day that by the time we master them they are of no use to us. Youth is the time to study wisdom, age the time to practice it. Experience is always instructive, I admit, but it is only useful in the time we have left to live. When death is already at the door, is it worth learning how we should have lived?
\flright{\textsc{Jean-Jacques Rousseau}}
\end{quotex}


Toward the end of his life, \textbf{Jean-Jacques Rousseau} published a series of ten meditations published as \emph{Reveries of the Solitary Walker}. Aside from his outbursts of paranoia, megalomania, and bitterness, they are instructive in the exploration of the spontaneous flow of thoughts. Like Heidegger, Rousseau walked in the woods. Julius Evola and Guido De Giorgio preferred the mountains. I walk on the beach with the same effect. My doctor wants me to do more of a ``power walk'', but that would be forcing myself into an external standard and it would impede the free flow of ideas.

First, we must consider what it means to be ``solitary''. Despite his writings on the social contract and the ``general will'', in the end he becomes a solitary, and we are left to wonder how such a one fits into conventional society. \textbf{Rene Guenon} was also a solitary contemplative; he wrote about Tradition and the hierarchical organization of a traditional society, but, for him, he allowed himself to pick and choose which Tradition to adhere to. His recommendation was that a man should follow the dominant Tradition of his city. Although not at all like the revolutionaries who flout every social convention, there is still a certain uneasiness about that. The fundamental question revolves around which traditions are natural and organic versus those that are merely conventional, appropriate only for a certain time and place.

The best way to read Rousseau's reveries is to regard them as a struggle to understand certain esoteric ideas without proper training or doctrine to follow. He recognizes the need for renunciation, not only for worldly vanities, but also for arbitrary opinions in his mind that may have arisen from emotional needs. He recognizes a cosmic moral law, despite how some interpret his earlier philosophical efforts. But most important is the process of ``letting go'', \emph{gelassenheit}, allowing thoughts to arise without controlling them. Ultimately, one must learn to live by one's intuition, as though every event is directed from a transcendent source. For complete liberation, the egoic idea that I am a ``thinker'' and a ``doer'', must be released in favour of a ``letting be''.

\paragraph{The Art of Dying}


These are things you need to start thinking about while you are young. The answers you come up with determine how you should live, but an old man needs to be more concerned about how he should die. Rousseau continues:

\begin{quotex}
We enter the race when we are born and we leave it when we die. Why learn to drive your chariot better when you are close to the finishing post? All you have to consider then is how to make your exit. If an old man has something to learn, it is the art of dying, and this is precisely what occupies people least at my age; we think of anything rather than that. Old men are all more attached to life than children, and they leave it with a worse grace than the young. This is because all their labours have had this life in view, and at the end they see that is has all been in vain. When they go, they leave everything behind, all their concerns, all their goods, and the fruits of all their tireless endeavours. They have not thought to acquire anything during their lives that they could take with them when they die.
\end{quotex}


Saints live in the vivid awareness of their necessary death. So, too, do warriors, for whom any day of battle may be their last. Living in an area populated with retirees, I see daily what Rousseau noted. There is an obsession with being ``youthful'', and they live as though the times of their youth are eternal. Recently, a college buddy called and asked me if I had seen some TV special about the Beatles. Aside from the metaphysical question of whether they were really the ``Beatles'' if half the original members were missing, I had no interest in the TV special. But that is not unusual, most ``boomers'' hold the same worldview and listen to the same music as though they were still 19 years old.

A sexagenarian told me that going for a hot air balloon ride was on her ``bucket list'', i.e., things she wanted to do before she died. My desire is to become free of mortal sin; she is much more sanguine about achieving her goal than I am.

\paragraph{Philosophers for Young Men}


I don't understand the fascination with Nietzsche in some circles. He was a beta male, a ``flop with chicks''. There is the famous story of his unrequited love, but, unlike a Dante, he derived no inspiration from it. It is curious that he is becoming popular among women, which is his fate despite his opinions of them. Apparently, they don't believe that Nietzsche was referring to them, although he actually was.

Nietzsche adapted Schopenhauer's philosophy by flattening it out and naturalizing it. Schopenhauer and Rousseau, despite their flaws, are more suitable for young men. They actually lived a life, filled with love and strife, and it shows up in their philosophizing. Moreover, they were brilliant stylists. Unfortunately, I do not know German well enough to read Schopenhauer in the original, but Rousseau's French is mesmerizing. Schopenhauer recognized kindred spirits in the rishis of the Upanishads, in the Buddha, and in the Neoplatonic mysticism of the Church Fathers. His philosophical system would be true only if the Noumenon, or the Will, was not a Subject, not an ``I'' or if the Father did not beget the Logos. But we will get back to that soon in an explication of Solovyov's metaphysics.

\paragraph{Exoteric Philosophy}


Rousseau understood that philosophy needed to arise from lived experience. Perhaps if he had been exposed to an esoteric tradition, he would have had better grounding for his own intellectual development. Nevertheless, it is instructive to see his own struggles with this idea.

\begin{quotex}
Thrown into the whirlpool of life while still a child, I learned from early experience that I was not made for this world, and that in it I would never attain the state to which my heart aspired. Ceasing therefore to seek among men the happiness which I felt I could never find there, my ardent imagination learned to leap over the boundaries of a life which was as yet hardly begun. ... This desire has at all time led me to seek after the nature and purpose of my being with greater interest and determination that I have seen in anyone else. I have met many men who were more learned in their philosophizing, but their philosophy remained as it were external to them. ... They studied human nature in order to speak knowledgeably about it, not in order to know themselves; their efforts were directed to the instruction of others and not to their own inner enlightenment.
\end{quotex}


This is certainly the way of academics who aspire to erudition. There was an instructive example in the comments recently about someone and his interlocutors who criticized Tradition on the grounds that it left out the ``saving grace of Christ''. They are forgetting, of course, that no one can quite agree on how that works. On the other hand, the esoteric understanding is precise: by calming the perturbations of the soul, the spirit can impregnate the mind and give birth to the Logos. That makes no sense from the exoteric point of view because it cannot be discussed by the discursive mind.

\paragraph{Fruits of Meditation}


Although Rousseau abandoned Catholicism for Calvinism, the style of meditation he practiced seems to derive from that spiritual movement that came to be called ``quietism.''

\begin{quotex}
The rural solitude in which I spent the best days of my youth, and the reading of good books which completely absorbed me, strengthened my naturally affectionate tendencies in her Madame de Warens company and led me to an almost \textbf{Fenelon}-like devotion. Lonely meditation, the study of nature and the contemplation of the universe lead the solitary to aspire continually to the maker of all things and to seek with a pleasing disquiet for the purpose of all he sees and the cause of all he feels. When my destiny cast me back into the torrent of this world I found nothing there which could satisfy my heart for a single moment.
\end{quotex}


Yet, up until the age of forty, he found himself ``living at random with no rational principles.'' At that point, he gave up many of the vanities that had been important to him. He also spontaneously discovered the necessity for the First Trial:

\begin{quotex}
I did not confine my reformation to outward things. Indeed I became aware that this change called for a revision of my opinions, which although undoubtedly more painful was also more necessary ... A great change which had recently come over me, a new moral vision of the world which had opened before me, the foolish judgments of men, whose absurdity I was beginning to sense without foreseeing how I was to fall victim to them.
\end{quotex}


Just as happens to many today, in his time Rousseau was shaken by ``\emph{certain modern philosophers who had little in common with the philosophers of antiquity}.'' He called them the ``\emph{missionaries of atheism, these overbearing dogmatists, who could not patiently endure that anyone should think differently from them on any subject whatsoever}.''

Budding esoterists and Hermetists should be careful about engaging in debate with those types of closed minded men. Until your understanding becomes firm and settled, they can do nothing for you except cause you doubt and uncertainty.

\paragraph{The Tendency to Self-Deception}


When clearing the mind of all opinions, there is the danger of self-deception. Reason is often helpless in the face of strong desires. Rousseau writes:

\begin{quotex}
It is hard to prevent oneself from believing what one so keenly desires, and who can doubt that the interest we have in admitting or denying the reality of the Judgment to come determines the faith of most men in accordance with their hopes and fears. All this may have led my reason astray ... But what I feared most was to endanger the eternal fate of my soul for the sake of those worldly pleasures which have never seemed very precious to me.
\end{quotex}


Ultimately, he claimed, he succeeded in freeing up his mind and was able to resist specious arguments and insoluble objections. He wrote:

\begin{quotex}
In this untroubled state of mind I find not only self-contentment, but the hope and consolation which my situation requires. A solitude so complete, so permanent and in itself so melancholy, the ever-present and constantly active animosity of all the present generation, the humiliations which they constantly heap on me, all this inevitably depresses me from time to time; uncertainty and worrying doubts still return occasionally to trouble my soul and fill it with gloom. Then I feel the need to recall my former conclusions ... Thus I reject all new ideas as fatal errors which have only a specious appearance of truth and are only fit to disturb my peace of mind.
\end{quotex}


\paragraph{Leaving Life}


Finally, Rousseau admits that it is the virtuous life that matters and not the attraction to the glamour of the world. He concludes this meditation:

\begin{quotex}
Patience, kindness, resignation, integrity and impartial justice are goods that we can take with us and that we can accumulate continually without fear that death itself can rob us of their value. It is to this one useful study that I devote what remains of my old age. And I shall be happy if by my own self-improvement I learn to leave life, not better, for that is impossible, but more virtuous than when I entered it.
\end{quotex}



\flrightit{Posted on 2014-03-10 by Cologero}

\begin{center}* * *\end{center}

\begin{footnotesize}
\begin{sffamily}

\texttt{JA on 2014-03-11 at 11:41 said:}

I think Nietzsche is beloved of the ADD generation in part because his style is to make his points in self contained brief aphorisms that -- in comparison to Schopenhauer, for example, are quite easy to read in one sitting.

White ole mad Freddy was indeed a failure with chicks, I doubt any man who followed a traditional life style to the letter would be a master of the boudoir... ..remember that Crowley seduced Coomaraswamy's wife... ..the black is more seductive than the white... ... 

\hfill

\texttt{Pickman on 2014-03-28 at 09:07 said:}

``Patience, kindness, resignation, integrity and impartial justice are goods that we can take with us and that we can accumulate continually without fear that death itself can rob us of their value.''

This wisdom is gold. I contemplate death and my departure all the time. Thanks for sharing. The longer we live we run the risk of accumulating baggage from our experiences with bad folk and degenerates. Unfortunately I wrestle with my daemons and scars in this case and wish to freed from lessons of harsh influences, ones that I can relate to Rousseau in this case. I was mean were these worth the cost at such late times beyond their use when we received them? I'm not so sure either. In any case I prefer to take a more than human transcendent Olympian approach, and this is what differentiates the Traditional man from mere petty worldly reactionaries of a heavy lower disposition.

Impulsive retards exist and we may become aware of their ``contaminating presence'' but let us keep true to ourselves above this foray. In better days a physical banishing of the wretched would have been possible.

Alas:

``The best revenge is to be unlike him who performed the injury.''

``The object of life is not to be on the side of the majority, but to escape finding oneself in the ranks of the insane.''

``Here is a rule to remember in future, when anything tempts you to feel bitter: not ``This is misfortune,'' but ``To bear this worthily is good fortune.''

``You are a little soul carrying about a corpse, as Epictetus used to say.''

-- MAA, Mediations (one could pull endlessly from this philosophy).

In essence, I find it interesting to read the experiences of older men, as vitality and facilities wane and cynicism bites harder, how in elder years they prepare themselves for their passing (extreme unction). Additional years spent may not be healthy for a man so leaving beautified is important to ponder. Basically, why waste ones time with petty BS, aim for greatness.

``Just that you do the right thing. The rest doesn't matter. Cold or warm. Tired or well-rested. Despised or honored. Dying... or busy with other assignments. Because dying, too, is one of our assignments in life. There as well: ``To do what needs doing.'' Look inward. Don't let the true nature of anything elude you. Before long, all existing things will be transformed, to rise like smoke (assuming all things become one), or be dispersed in fragments... to move from one unselfish act to another with God in mind. Only there, delight and stillness... when jarred, unavoidably, by circumstances, revert at once to yourself, and don't lose the rhythm more than you can help. You'll have a better grasp of the harmony if you keep going back to it.''

-- ibid,.

\end{sffamily}
\end{footnotesize}

